% Source: https://artofproblemsolving.com/community/c3h2313260_number_theory_a_lot_of_it
% Stable identity: aops-topic-2313260-q07
If $n=p_1^{e_1}p_2^{e_2}\cdots p_k^{e_k}$, then $n$ has $(e_1+1)(e_2+1)\cdots(e_k+1)$ positive divisors. For a given list of exponents, $n$ is smallest when the largest exponent goes to $2$, the next to $3$, and so on.

\textbf{(a)} Write $24$ as a product of factors $e_i+1\ge2$ in all ways and take the smallest number for each:
\[\begin{array}{ccccccc}
24&12\cdot2&8\cdot3&6\cdot4&6\cdot2\cdot2&4\cdot3\cdot2&3\cdot2\cdot2\cdot2\\\hline
2^{23}&2^{11}\cdot3&2^7\cdot3^2&2^5\cdot3^3&2^5\cdot3\cdot5&2^3\cdot3^2\cdot5&2^2\cdot3\cdot5\cdot7\\
8388608&6144&1152&864&480&360&420
\end{array}\]
The smallest is $\boxed{360}$.

\textbf{(b)} Each positive divisor $d$ gives the two divisors $d$ and $-d$, so we need exactly $12$ positive divisors. The products are $12$, $6\cdot2$, $4\cdot3$, $3\cdot2\cdot2$, giving $2^{11}=2048$, $2^5\cdot3=96$, $2^3\cdot3^2=72$ and $2^2\cdot3\cdot5=60$. The smallest is $\boxed{60}$.
